颅内动脉瘤是一种常见的脑血管异常，破裂后可引起蛛网膜下腔出血，危险性较高。CTA 能较清楚地显示载瘤血管和瘤体形态，但由于切片多、分叉复杂且存在骨伪影，人工阅片仍可能出现漏检或判断不一致。动脉瘤破裂风险又同时受到年龄、基础病、瘤体形态和位置等多种因素影响，单纯依靠传统量表或统计模型，往往难以全面刻画这些变量之间的非线性关系。因此，研究一种能够从影像中稳定分割动脉瘤，并结合临床信息进行风险分层的辅助方法，具有实际应用价值。\par

针对这一需求，本文在毕业设计阶段实现了两条可复现的技术路线。首先，在分割任务中采用三维 Attention U-Net，并在基础结构上进一步集成 ASPP 多尺度上下文模块与三维坐标注意力模块。基础模型保留了 U-Net 的多尺度跳跃连接，并在解码端加入注意力门控，对来自编码器的特征进行筛选，使网络在大量背景体素中更关注疑似瘤体区域；增强模型则在瓶颈层通过 ASPP 扩大感受野，并利用坐标注意力沿深度、高度和宽度方向重标定特征，以提升对细小动脉瘤及复杂血管走向的表达能力。训练时采用病例级流式加载和三维补丁采样，并结合重采样与强度归一化，以便在有限显存下处理大体积 CTA 数据；损失函数由二元交叉熵和 Dice 损失组合而成，用于同时兼顾分类准确性和区域重叠度。推理时，模型对整例数据分块预测，再将结果拼接回原始空间，便于后续使用。\par

其次，在破裂风险预测任务中，本文面向表格型临床和影像衍生特征，构建基于交叉注意力的深度表格模型：数值与类别字段分别映射到统一隐空间，经多层交叉注意力融合后由多层感知机输出破裂概率。数据处理包括缺失值填补、低频类别合并、特征筛选，以及面向类别不平衡的加权损失与验证集阈值优选。\par

在固定训练/验证/测试划分下，本文进一步将所提模型与随机森林、逻辑回归、支持向量机、$k$ 近邻及融合模型进行离线比较。结果表明，融合模型在测试集上取得最高 F1 值 0.7330 和 AUC 0.8568；本文的交叉注意力模型取得 F1 值 0.7266 和 AUC 0.8546，整体性能接近最优结果，并优于单独的逻辑回归、支持向量机和 $k$ 近邻方法。这说明基于注意力的特征交互机制能够有效建模异构风险因素之间的非线性关系，但在当前小样本条件下，传统集成模型仍表现出较强稳健性。\par

整体实现以配置驱动方式串联分割与风险两条流程，并提供简易交互界面，支撑从体数据读入到辅助分析演示的闭环。受数据脱敏与伦理审批进度限制，分割定量指标需结合具体实验设置解读；风险预测结果目前仍基于固定划分下的一次离线对比，尚不构成多中心外部验证结论。\par
